% Generated from locale/tr/content/intuitionistic-logic/semantics/introduction.tex; do not hand-edit.


\olfileid[tr]{int}{sem}{int}

\olsection{Giriş}

Hiçbir mantık bir semantik olmadan doyurucu biçimde betimlenemez; sezgici
mantık da bunun istisnası değildir. Klasik mantıkta !!{valuation}s{}i temel alan
semantik kanonik iken sezgici mantık için birbiriyle yarışan birkaç semantik
vardır. Bunların hiçbiri, bağlaçların anlamlarına sezgici bakımdan kabul
edilebilir bir açıklama getirme yönünden bütünüyle doyurucu değildir.

Kiplik mantıklarının semantiğine benzeyen ve !!{relational model}s{}i temel alan
semantik belki de bunların en yaygın olanıdır. Bu semantikte
!!{propositional variable}s dünyalara atanır ve bu dünyalar bir erişilebilirlik
bağıntısıyla ilişkilendirilir. Bu bağıntı her zaman bir kısmi sıralamadır; yani
yansımalı, ters simetrik ve geçişlidir.

Sezgisel olarak bu dünyaları bilgi durumları ya da ``kanıtsal durumlar'' diye
düşünebilirsiniz. Bir~$w'$ durumu~$w$'den ancak ve ancak bildiğimiz kadarıyla
$w'$ olası (gelecekteki) bir bilgi durumuysa, yani~$w$'de bilinenlerle
bağdaşıyorsa erişilebilirdir. Bir önerme bilindikten sonra yeniden bilinmez
hâle gelemez; yani $!A$,~$w$'de biliniyorsa ve~$Rww'$ ise $!A$,~$w'$'de de
bilinir. Dolayısıyla ``bilgi'', erişilebilirlik bağıntısına göre monotondur.

Epistemik mantıktaki gibi ``$!A$ biliniyor'' ifadesini ``bütün epistemik
alternatiflerde doğru'' diye tanımlarsak, $!A \land !B$,~$w$'de ancak bütün
epistemik alternatiflerde hem $!A$ hem de~$!B$ biliniyorsa bilinir. Ancak bilgi
monoton ve $R$ yansımalı olduğundan bu, $!A \land !B$'nin $w$'de bilinmesinin
ancak ve ancak $!A$ ile $!B$'nin~$w$'de bilinmesine eşdeğer olduğu anlamına
gelir. Aynı nedenle
$!A \lor !B$, $w$'de ancak ve ancak bunlardan en az biri biliniyorsa bilinir.
Dolayısıyla $\land$ ile $\lor$ için bağlaçların doğruluk koşulları klasik
mantıktakilerle örtüşür.

Ne var ki koşulun doğruluk koşulları klasik mantıktakilerden farklıdır. $!A
\lif !B$,~$w$'de ancak ve ancak hiçbir $w'$ için, $Rww'$ geçerliyken $!A$'nın
$!B$ de bilinmeden bilinmesi söz konusu değilse bilinir. Bu, $!A$'nın bilinmediği ya da $!B$'nin~$w$'de
bilindiği koşulla aynı değildir. Çünkü ne $!A$'yı ne de $!B$'yi~$w$'de
biliyorsak, gelecekteki bir~$w'$ epistemik durumu bulunabilir; bu durum için
$Rww'$ geçerlidir ve $w'$'de $!A$,~$!B$ de öğrenilmeksizin bilinebilir.

$\lnot !A$'yı ancak~$!A$'yı bildiğimiz olası hiçbir gelecek epistemik durum
yoksa biliriz. Buradaki düşünce şudur: $!A$ bilinebilir olsaydı, olası bir
gelecek epistemik durumda~$!A$ bilinir hâle gelirdi. $\lfalse$'ı
bilemeyeceğimizden, o gelecek epistemik durumda $!A$'yı bilir ama~$\lfalse$'ı
bilmezdik.

Bu yorumda üçüncü hâlin olanaksızlığı ilkesi geçersiz olur. Çünkü henüz
bilmediğimiz ama ileride bilebileceğimiz bazı $!A$'lar vardır. Böyle
!!a{formula}~$!A$ için hem $!A$ hem de~$\lnot !A$ bilinmez; dolayısıyla $!A
\lor \lnot !A$ da bilinmez. Fakat örneğin $\lnot(!A \land \lnot !A)$'yı
biliriz. Çünkü hem $!A$'yı hem de $\lnot !A$'yı bildiğimiz hiçbir gelecek durum
mümkün değildir; bunu~$!A$'yı ya da $\lnot !A$'yı bilip bilmediğimizden
bağımsız olarak biliriz.

Sezgici mantık için kullanılabilen tek semantik !!^{relational model}
değildir. Topolojik semantik de bir başkasıdır: burada önermeler bir topolojik
uzaydaki açık kümeler, bağlaçlar da bu kümeler üzerindeki işlemler olarak
yorumlanır (örneğin $\land$ kesişime karşılık gelir).

